\section{Finite realization by compatible intervals}
\label{sec:finite-realization}

This section lifts the finite cutoff chosen by the scalar argument back to the physical labelled layer.  Only after $X$ has been fixed do we use Haxell to select mutually compatible surviving rows.  The uniform simple-quotient supports are then lifted to covering affine translations by Corollary~\ref{cor:graded-affine-translation}, and joint compatibility strictifies their labelled composite.

Let $\eta>0$ be the reciprocal-mass margin in Proposition~\ref{prop:wide-start}.  By \eqref{eq:role-summable}, choose $B_{\mathrm{mass}}$ such that
\begin{equation}\label{eq:resource-tail}
\sum_{Q>B_{\mathrm{mass}}}\omega_Q<\eta.
\end{equation}
Choose also $B_{\mathrm{pack}}$ so large that every compatible subfamily $\mathcal A_Q(B)$ occurring after that point has cardinality at least the quota denominator $L_P$ of Proposition~\ref{prop:haxell-quota}.  This is possible by \eqref{eq:family-size} and Proposition~\ref{prop:wide-start}.  Apply Theorem~\ref{thm:scalar-terminality} with
\begin{equation}\label{eq:global-start-cutoff}
B_\dagger=\max\{B_{\mathrm{mass}},B_{\mathrm{pack}}\}.
\end{equation}
The relevant choices are made in the order
\begin{equation}\label{eq:realization-dependency}
\begin{aligned}
(B_{\mathrm{mass}},B_{\mathrm{pack}})&\longrightarrow B\longrightarrow X,\\
X&\longrightarrow P_X\longrightarrow
\{\mathscr L_X(Q):B<Q\le X\}\longrightarrow
(\mathscr T_1,\ldots,\mathscr T_m)\longrightarrow R_X.
\end{aligned}
\end{equation}
Here $B$ comes from Theorem~\ref{thm:scalar-terminality}, a finite $X$ is fixed before Haxell is invoked, $P_X$ is then chosen by Haxell, the uniform local theorem is applied to the surviving rows to form the physical label families $\mathscr L_X(Q)$ and covering translations $\mathscr T_j$, and joint compatibility finally strictifies their composite to $R_X$.

Fix $X\in\mathcal H_B$.  Only finitely many prime-power steps lie between $B$ and $X$.  For each such $Q$, Proposition~\ref{prop:wide-start} gives
\begin{equation}\label{eq:survivor-density}
|\mathcal A_Q(B)|
\ge\sigma|\mathcal A_Q|
\ge\kappa\sigma\frac Q{\log Q}.
\end{equation}
The classes $\mathcal A_Q(B)$ are pairwise disjoint: by the construction in Section~\ref{sec:signed-inverse}, a retained interval uniquely recovers its top prime-power stage $Q$.  Join two members of their finite union when they are incompatible as interval components.  Proposition~\ref{prop:signed-inverse-capabilities} gives a uniform maximum degree.  Proposition~\ref{prop:haxell-quota}, applied to this partition, yields an independent set $P_X$ with
\begin{equation}\label{eq:packed-density}
|P_X\cap\mathcal A_Q(B)|
\ge\rho_P|\mathcal A_Q(B)|
>\rho\frac Q{\log Q}.
\end{equation}
Hence
\[
P_X(Q):=P_X\cap\mathcal A_Q(B)
\in\mathcal D_\rho(Q)
\]
for every step through $X$.

\subsection{From the simple quotient to observable translation}

At a prime-power step $Q$, take
\[
H=F_{<Q},
\qquad
K=F_Q,
\qquad
K/H=S_Q.
\]
The general exact sequence \eqref{eq:affine-extension-exact} specializes to
\begin{equation}\label{eq:simple-extension}
\begin{CD}
0 @>>> S_Q @>>> A/F_{<Q} @>>> A/F_Q @>>> 0.
\end{CD}
\end{equation}
Thus Corollary~\ref{cor:graded-affine-translation} is precisely the lifting statement from a simple-quotient count support to a covering observable translation between the corresponding affine torsor fibres.

For each later step $B<Q=p^e\le X$, define
\begin{equation}\label{eq:physical-local-labels}
\mathscr L_X(Q)
=
\left\{
\bigcup_{I\in O}I:
O\subseteq P_X(Q),\ |O|\le p
\right\}.
\end{equation}
Because $P_X$ is independent, every displayed union belongs to $\mathcal C$.  By \eqref{eq:signed-label-filtration}, its residue lies in $F_Q$, and its image in $S_Q\times\N$ is exactly the corresponding point of
\[
\Gamma_{\phi_Q|_{P_X(Q)}}.
\]
Theorem~\ref{thm:uniform-profile} therefore supplies the quotient-and-count covering hypothesis of Corollary~\ref{cor:graded-affine-translation} with $\beta=0$.  Write
\[
B<Q_1<\cdots<Q_m\le X
\]
for the possibly empty list of processed prime-power steps.  Put
\[
T_0=\operatorname{Fib}_{E_B}(\gamma_B)\times I_B,
\]
and for $1\le j\le m$ let
\[
\gamma_j=\pi_{B,Q_j}(\gamma_B),
\qquad
T_j=\operatorname{Fib}_{E_{Q_j}}(\gamma_j)\times I_{Q_j}.
\]
The filtration and the scalar interval change only at prime-power steps, so $T_m=T_X$ when $m>0$, while for $m=0$ one has $T_0=T_X$.  For each $1\le j\le m$, Corollary~\ref{cor:graded-affine-translation} gives a covering observable translation relation
\begin{equation}\label{eq:physical-local-translation-chain}
\mathscr T_j:T_{j-1}\rightsquigarrow T_j.
\end{equation}

\subsection{Joint compatibility and exact composition}

\begin{lemma}[Joint compatibility]\label{lem:joint-compatibility}
Let $S_0$ be a label of the affine initial cover.  Let $Q_1,\ldots,Q_m\le X$ be distinct later prime-power steps, and choose $L_i\in\mathscr L_X(Q_i)$ for each $i$.  Then
\[
(S_0,L_1,\ldots,L_m)\in D_{m+1}.
\]
\end{lemma}

\begin{proof}
All interval components occurring in the $L_i$ belong to the single independent set $P_X$, so they are pairwise disjoint and nonadjacent.  Each lies in some $\mathcal A_{Q_i}(B)$ and is therefore compatible with every initial-cover label, in particular with $S_0$.  Their literal union with $S_0$ is consequently an element of $\mathcal C$.
\end{proof}

Let $P_n$ be the fibre product of the edge sets in the finite chain
\[
R_B,\mathscr T_1,\ldots,\mathscr T_m.
\]
Lemma~\ref{lem:joint-compatibility} gives the factorization
\begin{equation}\label{eq:realization-factorization}
\begin{CD}
P_n @>{\widetilde\ell}>> D_n @>{\iota}>> \mathcal C^n,
\end{CD}
\qquad
\ell=\iota\circ\widetilde\ell.
\end{equation}
Hence the comparison $2$-cell of \eqref{eq:oplax-2cell} is an identity by \eqref{eq:rel-strictification}.  Since $R_B$ and every $\mathscr T_j$ are covering, their ordinary relational composite is covering.  Strictification therefore gives the covering labelled composite
\begin{equation}\label{eq:physical-master-composite}
R_X:=\mathscr T_m\odot\cdots\odot\mathscr T_1\odot R_B:\{0_M\}\rightsquigarrow T_X,
\end{equation}
with
\begin{equation}\label{eq:observable-master-composite}
U(R_X)=U(\mathscr T_m)\circ\cdots\circ U(\mathscr T_1)\circ U(R_B).
\end{equation}
For $m=0$, both composites are understood to be $R_B$ and $U(R_B)$, respectively.  Here
\begin{equation}\label{eq:finite-response}
T_X:=\operatorname{Fib}_{E_X}(\gamma_X)\times I_X
\end{equation}
and $\gamma_X$ is the image of $\gamma_B$ under the endpoint quotient map.

\subsection{Reciprocal-mass control}

If $S_0$ is a label in the affine initial cover, Proposition~\ref{prop:wide-start} gives $W(S_0)<2-\eta$.  Every local label at step $Q=p^e$ is a compatible union of at most $p$ members of $\mathcal A_Q$, so Proposition~\ref{prop:signed-inverse-capabilities} bounds its reciprocal mass by $\omega_Q$.  By \eqref{eq:mass-subadditive} and \eqref{eq:resource-tail}, every interval configuration $S$ labelling an edge of $R_X$ satisfies
\begin{equation}\label{eq:physical-W-less-2}
W(S)
\le W(S_0)+\sum_{B<Q\le X}\omega_Q
<2.
\end{equation}

\begin{proposition}[Finite realization through a cutoff]\label{prop:finite-realization}
For every $X\in\mathcal H_B$ there is a covering observable labelled relation
\[
R_X:\{0_M\}\rightsquigarrow T_X
\]
with $T_X$ given by \eqref{eq:finite-response}.  Every edge label $S$ of $R_X$ satisfies $W(S)<2$.
\end{proposition}
